\section{Simulated annealing comparison: full protocol and additional sizes}
\label{app:sa_protocol}
\paragraph{Instances, objectives and initialization.}
For each of assignment and strip Set Partitioning, we use 40 frozen instances
at each $n=10,12,14$, with $N=n^2$ binary variables. Assignment uses its
row-major permutation encoding ($2n$ equality rows); the strip construction
and its $3n$ equalities are specified in \cref{app:partition_protocol}.
Each instance is evaluated separately with $C(x)=x^\top Qx$ and the 32-scenario
entropic objective. The symmetric quadratic arrays are formed by symmetrizing
matrices with independent uniform entries in $[-7,7]$. Within each $(n,i)$,
both families use the same frozen objective coefficient arrays and entropic
scale $s=\sqrt{3n}$, although their feasible spaces differ. Initial data are
400 independent uniform feasible draws, with replacement: random permutations
for assignment and exact counting for Set Partitioning. All four methods use
the same initial draw matrix within a problem, objective and budget.
The experiment comprises 5,760 runs: two families, two objectives, three sizes,
40 instances, three budgets and four methods.

\paragraph{Data-initialized EELS and full-indicator TT.}
Both start by deduplicating and evaluating the initial data. At every outer
iteration EELS rebuilds from the current training set and enhances the graph
with $\gamma=1$. The full TT constructs the exact complete indicator once
inside the timed run, then copies its uniformly normalized amplitudes at each
iteration. Both variants train for one sweep, using learning rate 0.05,
charge degeneracy one, and normalized Boltzmann weights with temperature
equal to the current cost standard deviation (uniform weights at zero spread).
After training they draw up to 10,000 samples, using the legacy backend in
batches of 100 for deadline checks. The next training set contains at most
400 distinct lowest-cost points from the previous training set and the new
batch; no external feasible samples are injected. The incumbent records all
evaluated feasible points, including a partially completed final batch.
The number of iterations is time-limited rather than fixed. These deliberate
resets must not be conflated with the fixed-topology, persistent-amplitude
comparison in \cref{subsec:eels_optimization}, nor with ADD's external-data
selection experiment.

\paragraph{Informed SA.}
Optim.jl's \texttt{SimulatedAnnealing} runs ten equally timed chains from the
ten best distinct initial points. For assignment a proposal swaps the assigned
columns of two distinct, uniformly chosen rows. For Set Partitioning a fixed
patch is chosen uniformly from the tile supports. If an occupied tile crosses
its boundary, the proposal is a self-loop; otherwise the current patch cover
is replaced by a uniformly chosen exact cover of that patch. Patch covers
are precomputed. These proposals are symmetric and preserve feasibility.
Each chain cools geometrically in wall time from the standard deviation of
initial costs to $10^{-3}$ times that value, with a machine-epsilon floor.
The moves and patch catalogue provide problem-specific information beyond
the explicit constraints and initial examples given to data-initialized EELS.

\paragraph{Penalty SA.}
The same Optim implementation and ten starting points minimize
$C(x)+\alpha\|Ax-b\|_2^2$ on the whole binary cube, proposing one uniformly
chosen bit flip. There is no repair or feasibility-preserving move catalogue.
With $R$ and $\sigma$ the range and sample SD of initial costs, set
\begin{align}
s_C&=\max\{R,\sigma,10^{-8}\max(1,\max_{x\in\mathcal T_0}|C(x)|)\},
&\alpha&=10s_C,\\
v&=\operatorname{median}_{j:\|A_j\|^2>0}\|A_j\|^2,
&T_0&=\max\{\sigma,\alpha v/\log2\},\\
&&T_{\mathrm{end}}&=\max\{10^{-3}\sigma,10^{-10}s_C\}.
\end{align}
Geometric cooling is restarted within each chain. Ignoring the objective
change, $T_0$ gives a typical constraint-violating one-bit proposal approximately
one-half acceptance probability, rather than freezing every move at a feasible
start. The multiplier is fixed across all instances and objectives; this is not
an optimized penalty or a proof of global exactness. Only feasible raw costs
update the reported incumbent. Every objective evaluation on an infeasible
state nevertheless consumes a call and time. Residual computation, penalty
setup, feasible/infeasible counters and best penalized states are recorded.
No infeasible evaluated point had a penalized value below the initial feasible
best, but this empirical observation is not a global bound.

\paragraph{Timing, counters and reproducibility.}
Separate 10, 50 and 200-second runs reuse fixed per-instance seeds; they are not
prefixes of a single trajectory or independent seed replicates. The source
seed plus $10^6$ initializes SA's default RNG, plus $2\times10^6$ its proposal
RNG, and plus $3\times10^6$ the TT RNG. Paired data and input hashes are checked.
The cluster used Intel Xeon Gold 6226R CPUs at 2.90 GHz, Julia 1.12.5 and
Optim.jl 2.3.2. Each solver process has one Julia compute thread and one BLAS
thread; 32 independent processes share each exclusive allocation, with four
allocations at a time. No GPU is used. Timed work includes initial objective
calls, TT construction, resets, training, sampling, penalty evaluation and
search. Loading, JIT warmup, pre-run validation and informed-SA patch-catalogue
construction are excluded; this last exclusion favors informed SA and limits
end-to-end runtime interpretation.

Termination is cooperative. Improvements after the requested deadline are
excluded, while actual elapsed times are retained. In particular, for quadratic
assignment at $n=14$ with a nominal 10-second budget, full TT averages 22.8
seconds (maximum 27.7) and completes no iteration on any task. At 200 seconds
all methods perform search. Exact objective counters include initial calls,
duplicates, repeated Optim evaluations, self-loops and penalty-SA infeasible
points; post-run validation calls are excluded. The same evaluator is used
by all methods: quadratic costs sum the active submatrix of $Q$, while the
entropic cost uses stable log-mean-exp. Neither method receives incremental
cost differences or cached objective answers. Feasible-call counts are not
numbers of distinct feasible solutions. Memory allocation totals are not
peak live memory; process MaxRSS spans sequential tasks and is not used for
per-method memory comparisons.

\paragraph{Aggregation and inference.}
For each panel, average the per-task gain
$100(C_0-C_{\mathrm{best}})/|C_0|$ equally over its 40 instances. The curve
uses the exact union of improvement events with stepwise interpolation.
Pointwise asymmetric 95\% percentile-bootstrap intervals resample the 40
tasks 2,000 times (seed 20261014), using the same weights across time and
methods. They are neither simultaneous confidence bands nor task SDs.
For the objective-call axis, each method stops at the minimum final count
over all 40 tasks; endpoints from slower tasks are not extrapolated. The
100,000-call diagnostic is a post hoc slice of the 200-second trajectories,
not a separately tuned fixed-query test. Full assignment at $n=14$ does not
reach that slice on all tasks. Percentages are rounded to one decimal in
figures/tables, while machine-readable files retain full precision.

Paired tests use one-sided Wilcoxon signed ranks of raw baseline-minus-TT
cost differences; differences below $10^{-9}+10^{-11}\max(|C_a|,|C_b|)$
are numerical ties. Tests versus penalty SA form a Holm family of 72 comparisons
(36 groups times two TT variants). The original informed-SA analysis, extended
to include penalty SA as a third candidate, uses 108 comparisons in a separate
Holm family. Both are exploratory analyses of this benchmark, with one seed
per task. Signed-rank evidence is not a test specifically of mean percentage
gain. No global optimum, distributional equivalence, or universal solver
ranking is certified. The inexpensive synthetic objectives and exact-cover
structure limit extrapolation to expensive or qualitatively different tasks.

% SA_ENDPOINT_TABLE

% SA_PENALTY_TEST_TABLE

% SA_APPENDIX_FIGURES

